\documentclass[11pt]{article}

\usepackage[margin=0.85in]{geometry}
\usepackage{amsmath,amssymb}
\usepackage{booktabs}
\usepackage{enumitem}
\usepackage[hidelinks]{hyperref}

\setlist[itemize]{leftmargin=1.4em}
\setlist[enumerate]{leftmargin=1.6em}

\title{SYSC 5108 Exam Study\\Assignment 3, Question 4}
\author{Jesse Levine}
\date{}

\begin{document}
\maketitle

\section*{Question Summary}

The network has one sensing neuron, three hidden processing neurons, and one
processing output neuron:
\[
    1 \rightarrow 2 \rightarrow \{3,4\} \rightarrow 5.
\]
All processing neurons, including the output neuron, use ReLU activation. The
input is
\[
    x_1 = 0.5
\]
and the desired output is
\[
    t = 0.9.
\]
The question asks for the forward output, the backpropagation \(\delta\) values,
the gradients with respect to all weights, updated weights for learning rate
\(\alpha=0.1\), and the reduction in error after the update.

\section*{Course and Textbook Cross-References}

\begin{itemize}
    \item \textbf{Chapter 4 slides, Forward Propagation Generalization:} each
    layer computes a linear/affine weighted sum and then applies an activation.
    \item \textbf{Chapter 4 slides, Hidden Units:} ReLU is
    \[
        \phi(z)=\max(0,z).
    \]
    \item \textbf{Chapter 4 slides, Chain Rule and Backward Propagation:}
    backpropagation recursively computes gradients using the chain rule.
    \item \textbf{Goodfellow, Bengio, and Courville, \emph{Deep Learning},
    Chapter 6:} feedforward networks compute affine transformations followed by
    activation functions; ReLU units use \(g(z)=\max(0,z)\), and gradients are
    propagated backward through the computation graph.
\end{itemize}

\section*{Given Network Parameters}

From the diagram:
\[
\begin{aligned}
    w_{2,1}&=0.3, & w_{2,0}&=0.1, \\
    w_{3,2}&=0.2, & w_{3,0}&=0.1, \\
    w_{4,2}&=0.5, & w_{4,0}&=0.1, \\
    w_{5,3}&=0.4, & w_{5,4}&=0.6, & w_{5,0}&=0.1.
\end{aligned}
\]

Use squared-error loss with the usual \(\tfrac12\) factor:
\[
    E = \frac12(t-\hat{y})^2.
\]
The ReLU derivative is
\[
    \phi'(z)
    =
    \begin{cases}
        1, & z>0,\\
        0, & z<0.
    \end{cases}
\]
All weighted sums in this problem are positive, so every ReLU derivative used in
backpropagation is \(1\).

\section*{Part (a): Forward Pass}

\subsection*{Neuron 2}
\[
\begin{aligned}
    z_2 &= w_{2,0}+w_{2,1}x_1 \\
        &= 0.1+(0.3)(0.5) \\
        &= 0.25.
\end{aligned}
\]
Because \(z_2>0\),
\[
    a_2=\phi(z_2)=0.25.
\]

\subsection*{Neuron 3}
\[
\begin{aligned}
    z_3 &= w_{3,0}+w_{3,2}a_2 \\
        &= 0.1+(0.2)(0.25) \\
        &= 0.15,
\end{aligned}
\]
so
\[
    a_3=\phi(z_3)=0.15.
\]

\subsection*{Neuron 4}
\[
\begin{aligned}
    z_4 &= w_{4,0}+w_{4,2}a_2 \\
        &= 0.1+(0.5)(0.25) \\
        &= 0.225,
\end{aligned}
\]
so
\[
    a_4=\phi(z_4)=0.225.
\]

\subsection*{Neuron 5}
\[
\begin{aligned}
    z_5
    &=
    w_{5,0}+w_{5,3}a_3+w_{5,4}a_4 \\
    &=
    0.1+(0.4)(0.15)+(0.6)(0.225) \\
    &=
    0.1+0.06+0.135 \\
    &=
    0.295.
\end{aligned}
\]
Therefore,
\[
    \hat{y}=a_5=\phi(z_5)=0.295.
\]

\[
    \boxed{\hat{y}=0.295}
\]

\section*{Part (b): Delta Values}

Use the convention
\[
    \delta_j=\frac{\partial E}{\partial z_j}.
\]

\subsection*{Output neuron}

For the output neuron,
\[
\begin{aligned}
    \delta_5
    &=
    \frac{\partial E}{\partial z_5} \\
    &=
    (\hat{y}-t)\phi'(z_5) \\
    &=
    (0.295-0.9)(1) \\
    &=
    -0.605.
\end{aligned}
\]

\subsection*{Hidden neurons}

For a hidden neuron \(j\),
\[
    \delta_j
    =
    \phi'(z_j)\sum_k w_{k,j}\delta_k,
\]
where the sum is over neurons \(k\) that receive input from neuron \(j\).

Neuron 4:
\[
    \delta_4
    =
    \phi'(z_4)w_{5,4}\delta_5
    =
    (1)(0.6)(-0.605)
    =
    -0.363.
\]

Neuron 3:
\[
    \delta_3
    =
    \phi'(z_3)w_{5,3}\delta_5
    =
    (1)(0.4)(-0.605)
    =
    -0.242.
\]

Neuron 2:
\[
\begin{aligned}
    \delta_2
    &=
    \phi'(z_2)\left(w_{3,2}\delta_3+w_{4,2}\delta_4\right) \\
    &=
    (1)\left[(0.2)(-0.242)+(0.5)(-0.363)\right] \\
    &=
    -0.0484-0.1815 \\
    &=
    -0.2299.
\end{aligned}
\]

\[
    \boxed{
    \delta_5=-0.605,\quad
    \delta_4=-0.363,\quad
    \delta_3=-0.242,\quad
    \delta_2=-0.2299
    }
\]

\section*{Part (c): Gradients}

For a weight \(w_{j,i}\) feeding activation \(a_i\) into neuron \(j\),
\[
    \frac{\partial E}{\partial w_{j,i}}
    =
    \delta_j a_i.
\]
For a bias weight \(w_{j,0}\), the input is \(1\), so
\[
    \frac{\partial E}{\partial w_{j,0}}
    =
    \delta_j.
\]

\begin{center}
\begin{tabular}{lrl}
\toprule
Weight & Gradient & Calculation \\
\midrule
\(w_{5,3}\) & \(-0.09075\) & \((-0.605)(0.15)\) \\
\(w_{5,4}\) & \(-0.136125\) & \((-0.605)(0.225)\) \\
\(w_{5,0}\) & \(-0.605\) & \(\delta_5\) \\
\(w_{4,2}\) & \(-0.09075\) & \((-0.363)(0.25)\) \\
\(w_{4,0}\) & \(-0.363\) & \(\delta_4\) \\
\(w_{3,2}\) & \(-0.0605\) & \((-0.242)(0.25)\) \\
\(w_{3,0}\) & \(-0.242\) & \(\delta_3\) \\
\(w_{2,1}\) & \(-0.11495\) & \((-0.2299)(0.5)\) \\
\(w_{2,0}\) & \(-0.2299\) & \(\delta_2\) \\
\bottomrule
\end{tabular}
\end{center}

\section*{Part (d): Updated Weights}

Use gradient descent:
\[
    w_{\text{new}}
    =
    w_{\text{old}}-\alpha\frac{\partial E}{\partial w},
    \qquad \alpha=0.1.
\]

\begin{center}
\begin{tabular}{lrrr}
\toprule
Weight & Old value & Gradient & New value \\
\midrule
\(w_{2,1}\) & 0.300000 & -0.114950 & 0.311495 \\
\(w_{2,0}\) & 0.100000 & -0.229900 & 0.122990 \\
\(w_{3,2}\) & 0.200000 & -0.060500 & 0.206050 \\
\(w_{3,0}\) & 0.100000 & -0.242000 & 0.124200 \\
\(w_{4,2}\) & 0.500000 & -0.090750 & 0.509075 \\
\(w_{4,0}\) & 0.100000 & -0.363000 & 0.136300 \\
\(w_{5,3}\) & 0.400000 & -0.090750 & 0.409075 \\
\(w_{5,4}\) & 0.600000 & -0.136125 & 0.613613 \\
\(w_{5,0}\) & 0.100000 & -0.605000 & 0.160500 \\
\bottomrule
\end{tabular}
\end{center}

\section*{Part (e): Error Reduction}

The original error is
\[
\begin{aligned}
    E_{\text{old}}
    &=
    \frac12(t-\hat{y})^2 \\
    &=
    \frac12(0.9-0.295)^2 \\
    &=
    \frac12(0.605)^2 \\
    &=
    0.1830125.
\end{aligned}
\]

Now do a new forward pass using the updated weights.

Neuron 2:
\[
    z_2^{\text{new}}
    =
    0.12299+(0.311495)(0.5)
    =
    0.2787375,
    \qquad
    a_2^{\text{new}}=0.2787375.
\]

Neuron 3:
\[
    z_3^{\text{new}}
    =
    0.1242+(0.20605)(0.2787375)
    \approx
    0.181634,
    \qquad
    a_3^{\text{new}}\approx 0.181634.
\]

Neuron 4:
\[
    z_4^{\text{new}}
    =
    0.1363+(0.509075)(0.2787375)
    \approx
    0.278198,
    \qquad
    a_4^{\text{new}}\approx 0.278198.
\]

Neuron 5:
\[
\begin{aligned}
    z_5^{\text{new}}
    &=
    0.1605+(0.409075)(0.181634)+(0.6136125)(0.278198) \\
    &\approx
    0.405508.
\end{aligned}
\]
Thus
\[
    \hat{y}_{\text{new}}\approx 0.405508.
\]

The new error is
\[
\begin{aligned}
    E_{\text{new}}
    &=
    \frac12(0.9-0.405508)^2 \\
    &\approx
    0.122261.
\end{aligned}
\]

Therefore, the reduction in error is
\[
\begin{aligned}
    E_{\text{old}}-E_{\text{new}}
    &=
    0.1830125-0.122261 \\
    &\approx
    0.060751.
\end{aligned}
\]

\[
    \boxed{\text{Error reduction}\approx 0.06075}
\]

\section*{Final Answer}

\[
\boxed{\hat{y}=0.295}
\]

\[
\boxed{
\delta_5=-0.605,\quad
\delta_4=-0.363,\quad
\delta_3=-0.242,\quad
\delta_2=-0.2299
}
\]

\[
\boxed{
\begin{array}{c|c}
\text{Weight} & \text{Updated value} \\
\hline
w_{2,1} & 0.311495 \\
w_{2,0} & 0.122990 \\
w_{3,2} & 0.206050 \\
w_{3,0} & 0.124200 \\
w_{4,2} & 0.509075 \\
w_{4,0} & 0.136300 \\
w_{5,3} & 0.409075 \\
w_{5,4} & 0.613613 \\
w_{5,0} & 0.160500
\end{array}
}
\]

\[
\boxed{
E_{\text{old}}=0.1830125,\quad
E_{\text{new}}\approx 0.122261,\quad
E_{\text{old}}-E_{\text{new}}\approx 0.06075
}
\]

\section*{Exam Notes}

\begin{itemize}
    \item Since every \(z_j>0\), every ReLU derivative in this problem is \(1\).
    That is why the backpropagation is numerically simple.
    \item The sign convention here is \(\delta_j=\partial E/\partial z_j\). With
    this convention, use \(w_{\text{new}}=w-\alpha\,\partial E/\partial w\).
    \item The output \(0.295\) is below the target \(0.9\), so the gradients are
    negative and gradient descent increases the relevant weights.
\end{itemize}

\end{document}
